\documentclass[12pt]{article}
\usepackage[UTF8]{ctex}
\usepackage[a4paper,margin=1in]{geometry}

\title{月光下的纸船}
\author{}
\date{}

\begin{document}

\maketitle

在一座靠近江边的小城里，住着一个叫林夏的孩子。她最喜欢的事情，是在傍晚坐在旧石桥边，听水声慢慢流过桥洞。别人都觉得那不过是一条普通的河，可林夏总觉得，河水里藏着谁也不知道的秘密。

一天傍晚，天边刚刚染上橘红色，林夏从书包里撕下一页写满算术题的纸，小心翼翼地折成了一只纸船。她把纸船放进水里，轻轻说道：“如果你真的知道秘密，就替我把愿望带到远方吧。”

纸船顺着水流慢慢漂走。奇怪的是，别的纸船总会很快湿透下沉，可这只纸船却像披着一层淡淡的银光，在河面上稳稳前行。林夏看得出神，忍不住沿着河岸一路追了过去。

夜色渐渐降临，月亮升起来，清亮的月光落在河面上，把整条河照得像一条会发光的丝带。纸船漂到城外一片安静的芦苇荡边，忽然停了下来。林夏刚靠近，就看见芦苇中亮起一点一点温暖的光，像是无数小小的星星落进了人间。

那些光原来是一群萤火虫。它们绕着纸船飞舞，最后在空中聚成了一行忽明忽暗的字：\textit{“勇敢的人，才能看见被遗忘的路。”}

林夏怔住了。就在那一瞬间，芦苇荡深处竟显出一条以前从未见过的小路。小路两旁开满了白色野花，尽头是一座小小的木屋。木屋的窗子透出昏黄灯光，像是在等她很久了。

她鼓起勇气走过去，轻轻推开门。屋里坐着一位白发老奶奶，正在火炉边缝补一张渔网。老奶奶抬头笑了，说：“你终于来了。我等那个愿意相信纸船的人，已经很多年了。”

林夏有些紧张，小声问：“您知道这条河的秘密吗？”

老奶奶点点头，从桌上拿起一只玻璃瓶。瓶子里装着许多折得很小很小的纸船，每一只上面都写着一个愿望。她说：“很多人小时候都会把愿望交给河流，可长大以后，他们就忘了。河流没有忘，所以把这些愿望都送到了我这里。”

林夏透过玻璃，看见那些纸船上写着各种各样的话：有人想再见一次远方的朋友，有人希望生病的妈妈快点好起来，也有人只是盼望明天不要下雨，好让自己能痛快地踢一场球。

“那我的愿望呢？”林夏问。

老奶奶温和地看着她：“你的愿望还在路上，因为你并不是把它丢给了河流，而是选择自己追了上来。真正重要的愿望，不该只靠水带走，还要靠自己一步一步去寻找。”

林夏低头看着那只发光的纸船，忽然明白了。她一直希望自己能去更远的地方，看更大的世界，却又总害怕离开熟悉的小城。原来纸船带她来到这里，不是为了替她完成愿望，而是为了告诉她，愿望本来就握在自己手里。

等她再抬起头时，木屋、火炉、老奶奶，还有漫天飞舞的萤火虫，都像被风吹散的月光一样，悄悄消失了。芦苇荡恢复了原来的安静，只有那只纸船停在岸边，船身上多了一行娟秀的小字：

\begin{center}
“朝着月亮升起的方向走，别害怕远方。”
\end{center}

第二天清晨，林夏把那只纸船夹进了自己的日记本。从那以后，她比从前更认真地读书，也开始把想去的地方、想做的事情一件件写下来。很多年后，当她真的离开小城，坐上开往远方的火车时，窗外正好有一轮明亮的月亮。

她忽然想起那个夜晚，笑着在心里轻声说道：“谢谢你，替我把秘密留下来。”

而那条小河仍旧安静地流过石桥，像是什么都没发生过。只是每到月圆的时候，城里的孩子们总会发现，河面上偶尔会漂过一只不会下沉的纸船，载着一个新的愿望，慢慢驶向看不见的远方。

\end{document}
